\documentclass[11pt]{article}
\usepackage[margin=1in]{geometry}
\usepackage{booktabs}
\usepackage{array}
\usepackage{xcolor}
\usepackage[colorlinks=true,linkcolor=blue,urlcolor=blue,citecolor=blue]{hyperref}

% Visible reminders: delete every \todo before submission.
\newcommand{\todo}[1]{\textcolor{red}{[TODO: #1]}}

\title{Can NSPARC Confocal Imaging Support Embryonic Cell-Lineage Tracking in
\textit{Caenorhabditis elegans}? A Pre-Specified Feasibility Benchmark Against Published diSPIM}
\author{\todo{Authors and affiliations}}
\date{\todo{Date}}

\begin{document}
\maketitle

\begin{abstract}
Light-sheet microscopy has been used to follow every nucleus in the developing
\textit{C.\ elegans} embryo for hours with no detectable phototoxicity. It is
unclear whether a laser-scanning confocal microscope with an image-scanning detector
(Nikon NSPARC) can do the same. We describe a pre-specified pilot that images
BV514 embryos (ubiquitous nuclear mCherry--histone) at 60~s and 180~s volume
intervals, alongside non-fluorescence-imaged controls, and measures manually
verified nuclear detection, lineage-link correctness, developmental timing and
hatching. Audit frames are selected per developmental stage before review, and all
analysis uses an open, tested pipeline that records the code version and inputs of
every report. The study is a feasibility benchmark, not an equivalence or safety trial.
In instrument rehearsals on non-study worms, a 41-plane stack (28.6~s) fitted the
60~s interval for 8~h with no missed volumes, while high-excitation single-plane
imaging caused rapid signal loss and loss of worm movement, supporting the use of
low excitation in the pilot.
\todo{Add one sentence of embryo results after the pilot.}
\end{abstract}

\section{Introduction}
Reconstructing the cell lineage of the \textit{C.\ elegans} embryo requires
three-dimensional imaging fast enough to follow dividing nuclei, gentle enough
not to disturb development, and sensitive enough to resolve crowded nuclei late in
embryogenesis. Duncan et al.\ imaged ten BV514 embryos (transgene \texttt{ujIs113},
mCherry--histone) on a dual-view inverted selective-plane illumination microscope
(diSPIM) at 20\textdegree C, one volume per minute for 7.5~h, and reported no
detectable phototoxicity based on developmental timing and hatching~\cite{duncan2019}.\footnote{The
main text of~\cite{duncan2019} reports 7.5~h and 451 volumes per embryo; a figure legend
in the same article refers to 6.4~h imaging sessions. We use the main-text values as the
cadence and duration target.}
Downstream, automated lineage tracing and identity assignment have been
benchmarked on light-sheet data~\cite{tran2026}.

The Nikon NSPARC detector is a 25-element array detector that records a small image
at each scanned point and is used for image-scanning microscopy; the vendor reports
improved signal collection and lower required excitation~\cite{nikon}. NSPARC
does not turn confocal illumination into a light sheet, and the vendor's phototoxicity
claims are not independent measurements. Whether this improvement is sufficient for
long, whole-embryo lineage tracking is an empirical question.

\textbf{Research question.} How far into embryogenesis can NSPARC confocal imaging
recover high-confidence nuclear identities and divisions at a useful temporal cadence
without altering normal development?

\section{Methods}
\subsection{Design and groups}
Embryos are randomly assigned to three groups, blocked by imaging day, operator,
strain and mounting batch: \textbf{NSPARC\_60s} (one whole-embryo Z stack every
60~s for 450~min, 451 volumes), \textbf{NSPARC\_180s} (every 180~s for 450~min,
151 volumes; a lower-dose arm) and \textbf{CONTROL} (same strain, mount and
temperature, without fluorescence excitation). Embryos in all groups are followed
for development and hatching on the same transmitted-light schedule, using the
lamp and not a laser-based transmitted channel, so that controls receive no laser
light. On this system the laser transmitted-light (TD) channel is illuminated by
the 561~nm line, which we confirmed during setup; it is therefore not used for
controls. \todo{Follow-up schedule.} A feasibility pilot of three embryos per group is
planned, expanding to about ten per group. All attempted acquisitions, failures
and exclusions are retained. The \emph{embryo} is the independent unit; frames and
nuclei are nested within embryos.

\subsection{Strain, mounting and environment}
Strain BV514 (\texttt{ujIs113}, ubiquitous mCherry--histone) is used where available.
Two- to four-cell embryos are mounted using a locally validated preparation
approved by the laboratory manager, at a target of 20\textdegree C with the measured
temperature recorded. Because unsealed mounts lost focus during a 15~h rehearsal
(Section~\ref{sec:rehearsal-results}), mounts for long runs are sealed against
evaporation and equilibrated on the stage before acquisition starts.
\todo{Mounting chamber, buffer, sealing method and equilibration time, once the local SOP is fixed.}

\subsection{Image acquisition}
Imaging uses a Nikon AX confocal on a Ti2 stand with the NSPARC detector, controlled
by NIS-Elements AR 6.20.02. Nuclei are imaged in a single channel: mCherry excited
at 561~nm and detected at 570--616~nm, with a Plan Apo $\lambda$D $60\times$ oil
objective \todo{NA; confirm that an oil objective is acceptable for live embryos,
as no $40\times$ water/silicone objective is installed}. Scanning is bidirectional
with no line integration, at a $512\times512$ scan format; zoom is set so that the
embryo fills the field at approximately 0.2~\textmu m per pixel (the
unzoomed pixel size at this format is 0.58~\textmu m). Z stacks use a 0.75~\textmu m
step over the full embryo depth (approximately 30~\textmu m, about 41 planes),
with top and bottom set on each embryo before every run; stored absolute Z
positions are cleared before each run, because a stored position from a previous
session moved the objective to the wrong height during setup. The laser shutter is
closed between volumes. NSPARC is operated in \todo{CF or SR mode; whether
summed-detector and pixel-reassignment reconstructions can both be produced from
the same acquisition}. \todo{Final zoom and Z range.}

At this format the software reports a frame time of about 0.49~s per plane. In
rehearsals at the same scan format ($10\times$ objective, Section~\ref{sec:rehearsal}),
a 41-plane stack took 28.3--28.7~s including Z motion (0.69~s per plane), leaving
about 31~s spare in each 60~s interval. The achieved stack time is measured again on
the first embryos with the $60\times$ objective and final zoom, and the 60~s arm is
not attempted if a stack takes 60~s or longer.

Focus is held with the Perfect Focus System (PFS) where it can lock on the sample;
the PFS state is recorded for every run. Without PFS, an unsealed slide drifted out
of focus within minutes even though the objective position was unchanged
(Section~\ref{sec:rehearsal-results}). \todo{PFS lock and offset on embryo mounts.}

Excitation is set to the lowest level that still gives usable nuclear detection,
and the power at the sample is measured in mW (laser percentage is not a dose).
For every embryo we record scan settings, measured power, achieved volume interval,
number of recorded volumes, temperature and the SHA-256 checksum of the raw ND2
file. Acquisition settings are saved as named templates per group so that every
embryo in a group is imaged identically. \todo{Measured power, achieved intervals.}

\subsection{Instrument rehearsals}\label{sec:rehearsal}
Before the pilot, the acquisition chain was rehearsed on adult and larval worms that
are not study samples: strain AML18 (\texttt{wtfIs3}, nuclear GFP and tagRFP in all
neurons) and other worms on unsealed slides, not immobilized, at room temperature,
with a Plan Apo $\lambda$D $10\times$ NA~0.45 objective. Two conditions were used:
\begin{itemize}
  \item \textbf{Low dose} (2026-10-23/24): 561~nm at 3\% laser, 570--616~nm, CF mode,
        $512\times512$ bidirectional, 41 planes at 0.75~\textmu m every 60~s, for 1~h
        and 8~h, PFS off.
  \item \textbf{High dose} (2026-11-10 and 11-12): 488~nm at 75.5\% and 561~nm at
        45.9--51.4\% laser plus the TD channel, $1024\times1024$, a single plane every
        4~s (2~h and 4~h) or 10~s (15~h), PFS off.
\end{itemize}
Laser power at the sample was not measured in the rehearsals, so the two conditions
are compared by settings only. Rehearsal data are used to check the acquisition
chain and to choose settings; they are not study data and do not enter any
outcome or feasibility target.

Each rehearsal file was analysed frame by frame from the raw ND2 with an
accompanying script (\texttt{scripts/analyze\_timelapse.py}; Z stacks are
max-projected). Timing was taken from the acquisition metadata. \emph{Signal} was
the mean of foreground pixels (above the frame median plus 6 median absolute
deviations) minus the median, expressed as a fraction of the first five frames.
\emph{Movement} was measured in two ways: the fraction of TD pixels that changed
by more than six robust standard deviations between consecutive frames (high-dose
runs only), and the turnover of the fluorescent foreground mask between frames
60~s apart ($|A \oplus B|/|A \cup B|$), which uses masks rather than intensities so
that bleaching is not counted as movement. Movement is reported as 30-min means
relative to the first 30~min of the same run. Image sharpness was the
variance of the Laplacian of the TD image, normalized by the squared mean intensity.

\subsection{Manual validation}
Audit frames are chosen before any results are reviewed. For each embryo, the
tracker's per-frame nucleus counts are used to select a fixed number of frames from
each stage ($<16$, 16--49, 50--99, 100--199, $\geq 200$ nuclei) with a recorded
random seed; the frame list is saved and its SHA-256 recorded before auditing, so
that the selection can be shown to precede the audit. Each audited frame is
examined slice by slice through the full depth, recording each nucleus's position
so that it is counted once, and compared with the tracker's detections to give
true positives (TP), false positives (FP) and false negatives (FN). Near-side and
far-side nuclei may be audited separately. As an internal check, TP\,+\,FP for each
frame must equal the tracker's count for that frame. Nuclei in division are counted
by a rule fixed before auditing \todo{rule}, and uncertain counts are flagged.
Division links are verified separately as correct or incorrect.
\todo{Number of frames per stage; number of reviewers.}

\subsection{Outcome measures}
\begin{itemize}
  \item Recall $=\mathrm{TP}/(\mathrm{TP}+\mathrm{FN})$, precision
        $=\mathrm{TP}/(\mathrm{TP}+\mathrm{FP})$ and
        $F_1=2\mathrm{TP}/(2\mathrm{TP}+\mathrm{FP}+\mathrm{FN})$, by stage and depth.
        The stage of a frame is taken from its total nucleus count.
  \item Fraction of verified parent--daughter links that are correct, by stage.
  \item Stage depth reached: the last stage with interpretable identities.
  \item Developmental timing relative to a common observed anchor (not raw movie time).
  \item Hatching as hatched/known outcomes, with unknown outcomes reported
        separately and never counted as failures.
  \item Achieved interval, incomplete volumes and measured excitation power.
\end{itemize}
Signal loss over time (photobleaching) is recorded as an acquisition quality measure.
It is not used as a measure of phototoxicity, which is assessed only through
developmental timing and hatching relative to CONTROL embryos.

\subsection{Provisional feasibility targets}
We pre-specify two internal targets, not literature-derived standards: recall of at
least 95\% on audited frames with at most 100 nuclei, and at least 90\% of verified
links correct. Targets are judged per group, with each embryo's own result reported
alongside. With no evaluable data a target is reported as ``not evaluable'', never
as a failure. Because a target computed only on early, easy frames can pass while
saying little about crowded stages, we report which stages contributed to each
result.

\subsection{Analysis software}
Analysis uses an open, dependency-free Python pipeline (\texttt{celegans-benchmark},
version 0.2.0) that validates the curated tables, rejects internally inconsistent
records, and reports descriptive statistics per embryo and group. Every report
includes a one-page summary and a provenance record with the software version, code
commit, settings and the SHA-256 of each input table. The pipeline is covered by
an automated test suite, including a study-sized fabricated dataset checked
against independent recomputation, and the same workflow was rehearsed on a public
\textit{C.\ elegans} embryo dataset~\cite{sugawara2021data} before the pilot; neither
the fabricated nor the public data contribute to results. Acquisition quality of each
raw time-lapse (interval, missed volumes, signal over time) is checked with the
pipeline's \texttt{run-check} command or the time-lapse script described in
Section~\ref{sec:rehearsal}.
\url{https://github.com/BioNanomics/celegans-nsparc-dispim-benchmark}
\todo{Release tag and commit hash used for the final analysis.}

\section{Results}
\subsection{Instrument rehearsals (non-study samples)}\label{sec:rehearsal-results}
Table~\ref{tab:rehearsal} summarizes the rehearsals. Acquisition timing was exact
in every run: no time point was late or missing. In the low-dose 8~h run, all 481
volumes were recorded at 60.00~s intervals (range 59.92--60.08~s) with a stable stack
time of 28.6~s, so the planned 60~s cadence with a 41-plane stack can be sustained for
the full study duration.

\begin{table}[h]
\centering
\footnotesize
\setlength{\tabcolsep}{4pt}
\caption{Instrument rehearsals on non-study worms ($10\times$ objective). Signal is the
background-subtracted foreground mean as a fraction of the start; movement is the last
30-min mean as a fraction of the first.}
\label{tab:rehearsal}
\begin{tabular}{@{}>{\raggedright\arraybackslash}p{1.0in}
                >{\raggedright\arraybackslash}p{1.75in}
                >{\raggedright\arraybackslash}p{0.65in}
                >{\raggedright\arraybackslash}p{1.45in}
                >{\raggedright\arraybackslash}p{1.05in}@{}}
\toprule
Run & Settings & Time points & Signal remaining & Movement, end vs start \\
\midrule
Low dose, 1~h  & 561~nm 3\%; 41 planes; every 60~s & 61/61
               & red 96\% at 1~h & 100\% (mask turnover) \\
Low dose, 8~h  & 561~nm 3\%; 41 planes; every 60~s & 481/481
               & red 93\% at 2~h, 65\% at 8~h & 105\% (mask turnover) \\
High dose, 2~h & 488~nm 75.5\%, 561~nm 51.4\%; 1 plane; every 4~s & 1801/1801
               & GFP 29\%, red 62\% at 2~h & 24\% (TD motion) \\
High dose, 4~h & as above & 3601/3601
               & GFP 10\%, red 26\% at 4~h & 4\% (TD motion) \\
\bottomrule
\end{tabular}
\end{table}

Signal fell much faster with high-dose single-plane imaging than with low-dose
volume imaging: GFP lost half its signal in 27--32~min, whereas the red signal in
the low-dose runs had not halved after 8~h despite 41 planes per time point. Worm
movement declined steeply under high-dose imaging (to 4\% of the initial value after
4~h in a field of crowded larvae, with an adult that stopped moving within about
1~h in the 2~h run), but not under low-dose imaging, where movement maps at 0, 4 and
8~h were similar. Because the two conditions also differed in slide, mount, sample
density and single-plane versus volume imaging, and no unexposed field was imaged on
the same slide, these observations point to excitation dose as a likely cause of the
loss of movement but do not establish it.

In a further 15~h high-dose rehearsal (488~nm 75.5\%, 561~nm 45.9\%, one plane every
10~s, 5401 time points, PFS off), the image lost focus progressively: TD sharpness
fell to 50\% of its starting value by 15~min and to about 26\% by 3~h, while the
recorded objective position stayed within 0.32~\textmu m. The sample, not the
objective, therefore moved relative to the focal plane, consistent with thermal
settling and drying of the unsealed mount. Fluorescence fell to background within
about 3~h. Movement could not be assessed after the first 1--2~h because defocus
masks it.

\subsection{Embryo pilot}
\todo{No embryo data yet. Report every embryo individually (group, settings, stage depth,
fate), then per-group summaries.}

\section{Discussion}
\todo{Write after results. Interpret as consistency with the published biological
benchmark, not as equivalence to diSPIM.}

The rehearsals already constrain the pilot. The 60~s cadence with a whole-depth
stack is achievable on this system. High excitation causes fast bleaching and is
associated with loss of worm movement, so the pilot uses the lowest usable 561~nm
excitation, measured in mW, and avoids the 488~nm line. Focus must be held by PFS,
a sealed and equilibrated mount, or a Z range wide enough to absorb drift. Bleaching
and phototoxicity are distinct: signal loss alone does not show harm, which is why
phototoxicity is judged only from development and hatching relative to CONTROL.

\subsection*{Limitations}
The pilot is not powered to show equivalence or safety. Comparison with Duncan
et al.\ is contextual, because their raw data are not matched to ours in strain
handling, optics, mounting or analysis. Our objective differs from the
water-immersion optics typical of light-sheet embryo imaging, and the stack time was
measured with a $10\times$ objective and must be confirmed with the $60\times$
objective and final zoom. Late-stage lineaging is limited by embryo movement and
nuclear crowding, and missing data there are not missing at random. Manual audit
effort rises steeply with nucleus number, so late-stage frames are sampled more
sparsely and their accuracy estimates are less certain. The rehearsals used
free-moving worms rather than embryos, unmeasured power at the sample, and
different slides for the low- and high-dose conditions without an unexposed control
field; their movement measures are descriptive and are confounded by worms leaving
the field or focal plane.

\begin{thebibliography}{9}
\bibitem{duncan2019} Duncan LH, Moyle MW, Shao L, et al. Isotropic light-sheet
microscopy and automated cell lineage analyses to catalogue \textit{Caenorhabditis
elegans} embryogenesis with subcellular resolution. \textit{J Vis Exp}.
2019;148:e59533. \url{https://doi.org/10.3791/59533}

\bibitem{tran2026} Tran MP, Peinado N, Zhao Z, Shah PK. Inference of
lineage-resolved cell identities in uncompressed \textit{C.\ elegans} embryos.
\textit{microPublication Biology}. 2026.
\url{https://pmc.ncbi.nlm.nih.gov/articles/PMC13242636/}

\bibitem{nikon} Nikon. AX / AX R with NSPARC (product information).
\url{https://www.microscope.healthcare.nikon.com/products/confocal-microscopes/ax/ax-ax-r-with-nsparc}

\bibitem{sugawara2021data} Sugawara K, \c{C}evrim \c{C}, Averof M. Supplementary data
(Fluo-N3DH-CE image and tracking data). Zenodo. 2021.
\url{https://doi.org/10.5281/zenodo.4630933}
\end{thebibliography}

\end{document}
